% from ega4-13-fr.tex
\oldpage[IV-3]{188}
Malheureusement la notion de morphisme équidimensionnel n'est pas stable par
changement de base~; c'est pourquoi dans de nombreuses questions il est plus
commode de travailler avec la notion de morphisme universellement ouvert, dont
l'étude fait l'objet des \S\S~14 et 15.

\subsection{Le théorème de semi-continuité de Chevalley}
\label{subsection:IV.13.1-fr}

\begin{lemma}[13.1.1]
\label{IV.13.1.1-fr}
Soient $Y$ un préschéma irréductible localement noethérien, $X$ un préschéma
irréductible, $f:X\to Y$ un morphisme dominant de type fini. Soit $\xi$
(resp. $\eta$) le point générique de $X$ (resp. $Y$) et soit
$e=\dim(f^{-1}(\eta))$. Alors, pour tout $x\in X$, toutes les composantes
irréductibles de $f^{-1}(f(x))$ sont de dimension $\geq e$.
\end{lemma}

\begin{proof}
La proposition est immédiate lorsque, pour tout $y\in f(X)$,
$\mathcal O_y$ est un \emph{anneau universellement caténaire}~: en effet, si
$x$ est point générique d'une composante irréductible $Z$ de $f^{-1}(y)$, il
résulte de (5.6.5), joint à (5.2.1), que l'on a
$e+\dim(\mathcal O_y)=\dim(Z)+\dim(\mathcal O_x)$~; mais en vertu de
($\mathbf 0$, 16.3.9), on a $\dim(\mathcal O_x)\leq\dim(\mathcal O_y)$,
d'où la conclusion dans ce cas.

Nous allons ramener le cas général à ce cas particulier. La question est
évidemment locale sur $Y$, et, compte tenu de (4.1.1.3), elle est aussi locale
sur $X$~; on peut donc se borner au cas où $Y=\operatorname{Spec}(A)$ et
$X=\operatorname{Spec}(B)$ sont affines et irréductibles, $B$ étant une
$A$-algèbre de type fini. En outre (1.5.4), on peut supposer $X$ et $Y$
réduits, donc $A$ et $B$ \emph{intègres} et, comme $f$ est dominant, $A$ est
alors un \emph{sous-anneau} de $B$. Considérons $A$ comme limite inductive de
ses sous-$\mathbf Z$-algèbres de type fini~; il résulte alors de (8.9.1) qu'il
existe une telle sous-algèbre $A_0$ et une $A_0$-algèbre de type fini $B_0$
telles que $B=B_0\otimes_{A_0}A$. Posons $Y_0=\operatorname{Spec}(A_0)$,
$X_0=\operatorname{Spec}(B_0)$, et soit $f_0:X_0\to Y_0$ le morphisme
correspondant à l'homomorphisme $A_0\to B_0$, de sorte que
$f=f_0\times 1_Y$. Il n'est pas évident \emph{a priori} que le préschéma $X_0$
soit intègre, mais nous allons voir qu'on peut se ramener à ce cas. Soit
$\eta_0$ le point générique de $Y_0$, de sorte que si $g$ est le morphisme
$Y\to Y_0$, on a $g(\eta)=\eta_0$~; par transitivité des fibres
(\textbf{I}, 3.6.4), on a
$f^{-1}(\eta)=f_0^{-1}(\eta_0)\otimes_{k(\eta_0)}k(\eta)$, et comme
$f^{-1}(\eta)$ est irréductible par hypothèse, il en est de même de
$f_0^{-1}(\eta_0)$ (4.4.1). Notre assertion résultera alors du lemme suivant~:

\begin{lemma}[13.1.2]
\label{IV.13.1.2-fr}
Soient $Y_0$, $Y$ deux préschémas intègres de points génériques $\eta_0$,
$\eta$, $g:Y\to Y_0$ un morphisme dominant, $f_0:X_0\to Y_0$ un morphisme
dominant tel que $f_0^{-1}(\eta_0)$ soit irréductible. Soit $X'_0$ l'unique
composante irréductible de $X_0$ rencontrant $f_0^{-1}(\eta_0)$
($\mathbf 0_{\mathrm I}$, 2.1.8), et désignons encore par $X'_0$ le
sous-préschéma fermé réduit de $X_0$ ayant $X'_0$ pour espace sous-jacent.
Supposons que le préschéma $X=X_0\times_{Y_0}Y$ soit intègre~; alors $X$ est
isomorphe à $X'_0\times_{Y_0}Y$.
\end{lemma}

En effet, si $j_0:X'_0\to X_0$ est l'injection canonique, qui est une immersion
fermée, $j=j_0\times 1_Y:X'_0\times_{Y_0}Y\to X_0\times_{Y_0}Y=X$ est une
immersion fermée. D'autre part, $X'_0$ contient la fibre
$f_0^{-1}(\eta_0)$, donc $X'_0\times_{Y_0}Y$ contient $f^{-1}(\eta)$~;
notons d'autre part que $f^{-1}(\eta)$ n'est pas vide (\textbf{I}, 3.4.7),
donc contient le point générique de $X$~; par suite l'image de $j$ est
nécessairement $X$ tout entier. Mais comme $X$ est intègre, le seul
sous-préschéma fermé de $X$ ayant $X$ pour espace sous-jacent est $X$
lui-même, donc $j$ est un isomorphisme.

Ce lemme étant établi, on peut donc supposer que $X_0$ est intègre~; pour tout
$y_0\in Y_0$, $\mathcal O_{y_0}$ est une $\mathbf Z$-algèbre essentiellement
de type fini, donc un anneau universellement

